(1) 高温等温理想气体 \(\Delta U=0\)，对外功
\[
W_{12}=\int_{V_1}^{V_2}\frac{nRT_h}{V}\,dV=nRT_h\ln r.
\]
因此吸热 \(Q_h=W_{12}=nRT_h\ln r>0\)。

(2) 可逆绝热满足 \(TV^{\gamma-1}=\mathrm{const}\)，所以 \(V_3/V_2=(T_h/T_c)^{1/(\gamma-1)}\)；回程 \(V_4/V_1=(T_h/T_c)^{1/(\gamma-1)}\)。相除可见低温等温压缩的体积比 \(V_3/V_4=V_2/V_1=r\)。

(3) 低温等温段对外功 \(W_{34}=nRT_c\ln(V_4/V_3)=-nRT_c\ln r\)，故排向冷源的热量大小 \(Q_c=nRT_c\ln r\)。两绝热段 Q=0、但其功与内能变化互相抵消；完整周期 \(\Delta U_{\rm cycle}=0\)，所以
\[
W_{\rm net}=Q_h-Q_c=nR(T_h-T_c)\ln r.
\]

(4) 热机效率
\[
\eta=\frac{W_{\rm net}}{Q_h}=1-\frac{T_c}{T_h}.
\]
由于可逆循环的熵收支 \(Q_h/T_h-Q_c/T_c=0\)，效率与 n、\(C_V\) 和 r 无关；若任一段不可逆，则熵产生正，实际效率低于该上限。
